% =============================================================
% Structured Representation for Code LLMs:
% When Format Can and Cannot Help a Bounded Learner, and Why
% =============================================================
% Theory companion to the empirical preprint. Each section is a
% separate file under sections/.
%
% Build with tectonic (self-contained, resolves natbib automatically):
%   cd paper-theory && tectonic main.tex

\documentclass[11pt]{article}

\usepackage[utf8]{inputenc}
\usepackage[T1]{fontenc}
\usepackage{lmodern}
\usepackage[margin=1in]{geometry}
\usepackage{microtype}
\usepackage{graphicx}
\usepackage{amsmath}
\usepackage{amssymb}
\usepackage{amsthm}
\usepackage{mathtools}        % \vcentcolon, \coloneqq helpers used in notation
\usepackage{bm}
\usepackage{booktabs}
\usepackage{array}            % ragged-right p-columns for the dense §10 ledger tables
\usepackage{xcolor}
\usepackage[round]{natbib}
\usepackage[hidelinks]{hyperref}
\usepackage[capitalize, noabbrev]{cleveref}

% ----- Theorem environments (numbered within section) -----
\theoremstyle{plain}
\newtheorem{theorem}{Theorem}[section]
\newtheorem{lemma}[theorem]{Lemma}
\newtheorem{proposition}[theorem]{Proposition}
\newtheorem{corollary}[theorem]{Corollary}
\theoremstyle{definition}
\newtheorem{definition}[theorem]{Definition}
\newtheorem{assumption}[theorem]{Assumption}
\newtheorem{example}[theorem]{Example}
\theoremstyle{remark}
\newtheorem{remark}[theorem]{Remark}
% ----- Falsifiable predictions (§8) -----
% Predictions are NOT proved results: they carry an independent counter (P8.x)
% so they are never confused with the theorems they derive from, and they live
% in remark style to keep the conjectural status visually distinct. Each
% prediction is stated with a sign/size claim and an explicit falsifier in its
% body.
\newtheorem{prediction}{Prediction}[section]
% cleveref auto-derives the singular name but not the plural for this custom
% environment; declare both so grouped/range references (e.g.
% \Cref{pred:capacity,pred:difficulty,pred:compute}) render "Predictions 8.2--8.4"
% rather than the literal "?? 8.2--8.4".
\crefname{prediction}{Prediction}{Predictions}
\Crefname{prediction}{Prediction}{Predictions}

% ----- Idealization-disclosure box -----
% Every invocation of the irrelevance theorem must state which idealization
% it relies on. This command standardizes that disclosure.
\newcommand{\idealization}[1]{%
  \par\vspace{2pt}\noindent\textit{Idealization disclosure.} #1\par\vspace{2pt}}

% =============================================================
% notation.tex — shared notation for the theory companion paper
% =============================================================
% Single source of truth for symbols across all sections. Every theorem,
% lemma, and proof draws its notation from here so that the four
% mathematical areas (measure-theoretic probability, information theory,
% statistical learning theory, and the two go/no-go layers) stay
% internally consistent. Edit here, not inline.
%
% Convention: areas are grouped and labelled. When a T1/T2 issue introduces
% a new symbol it adds it to the relevant group below with a one-line
% comment, rather than redefining \newcommand locally in a section.

% ----- Measure-theoretic probability (§2, T1-1) -----
\newcommand{\Xspace}{\mathcal{X}}            % input space (programs/representations source)
\newcommand{\Yspace}{\mathcal{Y}}            % label/target space
\newcommand{\Zspace}{\mathcal{Z}}            % representation codomain
\newcommand{\Bor}[1]{\mathcal{B}(#1)}        % Borel sigma-algebra of #1
\newcommand{\Prob}{\mathbb{P}}               % probability measure
\newcommand{\Expect}{\mathbb{E}}             % expectation
\newcommand{\given}{\,\vert\,}               % conditioning bar
\newcommand{\kernel}[1]{\mathsf{#1}}         % Markov kernel (predictor), e.g. \kernel{K}
\newcommand{\rep}{\phi}                      % representation map \phi : X -> Z
\newcommand{\risk}{\mathcal{R}}              % risk functional
\newcommand{\bayesrisk}{\risk^{\star}}       % Bayes risk
\newcommand{\loss}{\ell}                     % loss function

% ----- Information theory (§4, T1-2 / T2-1 / T2-2) -----
\newcommand{\MI}[2]{I(#1 ; #2)}              % mutual information
\newcommand{\Ent}[1]{H(#1)}                  % (differential) entropy
\newcommand{\CondEnt}[2]{H(#1 \given #2)}    % conditional entropy
\newcommand{\KL}[2]{D_{\mathrm{KL}}\!\left(#1 \,\Vert\, #2\right)} % KL divergence
\newcommand{\markov}{\rightarrow}            % Markov-chain link, X \markov \rep(X) \markov Y
\newcommand{\suff}{\;\succeq_{\mathrm{suff}}\;} % sufficiency preorder
\newcommand{\ratedist}{R(D)}                 % rate-distortion function

% ----- Statistical learning theory (§5, T1-3 / T2-4) -----
\newcommand{\hypclass}{\mathcal{H}}          % hypothesis class
\newcommand{\capac}[1]{\mathfrak{C}(#1)}     % capacity/complexity term (VC, Rademacher, ...)
\newcommand{\Rad}[1]{\mathfrak{R}_n(#1)}     % empirical Rademacher complexity at sample n
\newcommand{\emprisk}{\widehat{\risk}}       % empirical risk
\newcommand{\nsample}{n}                     % sample size
\newcommand{\advantage}{\Delta}              % representation advantage (bounded-regime gap)

% ----- Category theory (§3, T1-4 / T2-3) — GO/NO-GO -----
\newcommand{\catRep}{\mathbf{Rep}}           % category of representations
\newcommand{\ladder}[1]{\mathrm{C}#1}        % C-ladder object, e.g. \ladder{4} = C4
\newcommand{\ladderplus}{\mathrm{C1^{+}}}    % information-parity control C1+
\newcommand{\mor}{\xrightarrow}              % morphism arrow with label: \mor{f}
\newcommand{\retract}{\;\trianglelefteq\;}   % retraction relation (parity = retraction)
\newcommand{\roundtrip}[2]{#1 \!\to\! #2 \!\to\! #1'} % round-trip defect C4 -> C1+ -> C4'

% ----- Optimal transport + information geometry (§6, T1-5 / T2-5) — GO/NO-GO -----
\newcommand{\Wass}[1]{W_{#1}}                % Wasserstein-#1 distance
\newcommand{\preddist}[1]{\mu_{#1}}          % induced predictive distribution under condition #1
\newcommand{\Fisher}{\mathcal{I}_{F}}        % Fisher information
\newcommand{\paritydefect}{\delta_{\mathrm{par}}} % quantitative parity defect

% ----- Cross-cutting -----
\newcommand{\defeq}{\vcentcolon=}            % "defined as"
\newcommand{\indep}{\perp\!\!\!\perp}        % statistical independence


\title{Structured Representation for Code LLMs:\\
       When Format Can and Cannot Help a Bounded Learner, and Why}

\author{%
  \textbf{\large Takumi Shiraishi} \\[2pt]
  \href{mailto:takumi@wovol.com}{\texttt{takumi@wovol.com}}
}

\date{June 2026}

\begin{document}
\maketitle

\begin{abstract}
Does the format in which code is presented to a language model change what the
model can predict about it? We give a theoretical companion to the frozen
empirical record of \texttt{preprint-v1}, whose central pre-registered cell
returned a clean null. Making representations measurable maps and predictors
Markov kernels on a shared standard Borel space, we prove a two-halved answer.
\emph{Negatively}, under \emph{information parity} --- two representations
retaining the same mutual information with the target --- the data-processing
inequality closes at the Bayes optimum and the optimal risk is
format-invariant, so the recorded null is the \emph{predicted} information-limit
outcome rather than an anomaly; when parity fails, the inflation of Bayes risk
is controlled by the information lost (exactly under log loss, within a
square-root envelope otherwise). \emph{Positively} --- and never as a claim that
any effect exists --- a finite-capacity learner does not attain the information
limit, and the residual is exactly where format may act: at constant
information the achievable representation advantage is capped by a capacity
budget carrying no mutual-information term, two-sided in magnitude, and
non-vacuous only in a named finite-resource, capacity-separated regime. Two
further structural refinements show operational parity is a retraction whose
round-trip defect is invisible to both the mutual information and the Bayes
risk yet
dominates the swap gap of any fixed bounded predictor. We map the three frozen
Lumen outcomes to the regimes they occupy --- claiming consistency, not
confirmation --- and convert the localization into five falsifiable predictions
on the same apparatus, each with an explicit falsifier. The thesis: format
cannot supply information that is not there, so any genuine
code-representation effect on a bounded model is a capacity phenomenon of
bounded size, appearing only near the model's frontier.
\end{abstract}

\tableofcontents
\clearpage

% Section 1 — Introduction (charter §4.1; issue T4-1 prose, drawing on T1-2/T1-3)
\section{Introduction}
\label{sec:intro}

Does the \emph{format} in which a program is presented to a code language model
change what the model can predict about it? The question is usually asked
empirically --- pick a family of representations, hold the model fixed, and
measure the score --- and the empirical companion to this paper,
\texttt{preprint-v1}, asked it exactly that way on the Lumen apparatus: a frozen
model queried under a ladder of code representations C2--C4 against an
information-matched control C1${}^{+}$. Its central, pre-registered cell
returned a clean null (a rank-biserial effect $r_{\mathrm{rb}}=+0.10$ at raw
$p=0.36$, no rubric ceiling). This paper is the theoretical sibling of that
record. It does not revise the preprint, re-run its scorers, or touch its frozen
analysis; it asks \emph{why} the null is the outcome one should have expected,
and --- the harder half --- where an effect, if one exists at all, is permitted
to live.

The answer we defend has two halves, and the paper is honest about which half
each result belongs to. The first is negative and is a theorem. Once
``representation'' and ``predictor'' are made into measurable objects on a common
probability space (\cref{sec:prelim}), a representation is a measurable map that
can only \emph{coarsen} the input $\sigma$-algebra, never enrich it. The
data-processing inequality (\cref{lem:dpi}) then forbids any representation from
creating predictive information about the target, and at the Bayes optimum this
sharpens to an equality: under \emph{information parity} --- two representations
retaining the same mutual information with $Y$ --- the Bayes-optimal risk is
\emph{format-invariant} (\cref{thm:irrelevance}). At constant information the
optimum cannot tell the formats apart. The Lumen T1 null is, on this reading,
not a weak or disappointing positive but the \emph{predicted} outcome of the
information-limit regime, and treating it as anything else is the overclaim our
honesty guardrails exist to prevent.

The second half is positive, and it is deliberately \emph{not} a claim that any
effect exists. A frozen LLM is not the Bayes posterior; it is a
finite-capacity predictor that does not attain the information limit, and the
gap between what it achieves and that limit is precisely where format is free to
act. We bound that gap. For bounded learners the achievable
parity-preserving advantage is at most a capacity budget $\Phi$ built from an
approximation (expressivity) gap and a Rademacher estimation term, and carrying
\emph{no} mutual-information term whatsoever (\cref{thm:bounded-gap}); the bound
is one-sided and sign-agnostic --- it says how large an effect \emph{can} be,
never that one exists or which format wins (\cref{cor:size-gap}). Because the
budget is non-empty only in a named finite-resource, capacity-separated window
(\cref{prop:regime}), the theory does something a bare null cannot: it
\emph{localizes} the regime in which a representation effect is permitted,
forbids it everywhere else, and turns that localization into a list of
falsifiable predictions on the same apparatus (\cref{sec:predictions}).

\subsection{The format-versus-information question, stated formally}
\label{subsec:intro-question}

Informally the empirical question conflates two things that the theory must
separate: whether a representation carries different \emph{information} about the
target, and whether, at fixed information, a particular \emph{learner} extracts
different value from it. The Lumen C1${}^{+}$ control exists precisely to fix the
first so the second can be asked alone, and the formal apparatus of this paper
exists to say what ``fix the first'' means exactly. We write $\rep:\Xspace\to
\Zspace_{\rep}$ for a representation and ask of a pair $\rep,\rep'$:
\begin{enumerate}
  \item[(Q1)] \emph{Information.} Do $\rep$ and $\rep'$ retain the same
    predictive content, $\MI{Y}{\rep(X)}=\MI{Y}{\rep'(X)}$? This is
    \emph{information parity} (\cref{def:parity}); its categorical form is a
    retraction (\cref{sec:category}) and its quantitative failure is a transport
    defect (\cref{sec:parity-defect}).
  \item[(Q2)] \emph{Value to a bounded learner.} Given parity, can a
    finite-capacity predictor still do measurably better under one than the
    other? This is the representation advantage $\advantage(\rep,\rep')$
    (\cref{def:advantage}), and the bounded half of the paper is the statement
    that whatever value exists here is capacity, not information
    (\cref{sec:bounded}).
\end{enumerate}
Phrased this way the empirical null is the answer to a \emph{composite} of Q1 and
Q2: the C4-vs-C1${}^{+}$ comparison was engineered so Q1 is answered ``same''
(by construction), leaving the recorded effect to measure Q2 alone. The theorem
of \cref{sec:parity} is that the Q1 channel is genuinely closed at the optimum;
the bound of \cref{sec:bounded} is that the Q2 channel, while open, is narrow and
capacity-controlled. The negative half disposes of the information explanation;
the positive half says where to hunt for the capacity one.

\subsection{Contributions}
\label{subsec:intro-contributions}

\begin{enumerate}
  \item \textbf{A measure-theoretic frame in which the question is exact}
    (\cref{sec:prelim}). Representations become measurable maps and predictors
    become Markov kernels on a shared standard Borel space, so that ``the
    predictor only sees the representation'' is the statement that its kernel is
    $\repsig$-measurable, and the data-processing and sufficiency arguments
    become exact statements about conditional expectations rather than heuristics
    about ``information.''
  \item \textbf{An irrelevance theorem and its rate--distortion corollary}
    (\cref{sec:parity}). Under information parity the Bayes risk is
    format-invariant (\cref{thm:irrelevance}); when parity fails, the inflation
    of Bayes risk is controlled by the information lost --- exactly the lost bits
    under log loss, and a $\sqrt{\cdot}$ envelope for any bounded loss
    (\cref{cor:ratedist-gap}). This is the formal content of the negative half
    and the disciplined statement of guardrail~G1.
  \item \textbf{A bounded-learner advantage bound} (\cref{sec:bounded}). At
    constant information the achievable advantage is capped by a capacity budget
    $\Phi$ with no mutual-information term (\cref{thm:bounded-gap}), two-sided in
    magnitude (\cref{cor:size-gap}), and non-vacuous only in a named
    finite-resource window (\cref{prop:regime}). This is the positive half,
    stated as a bound and never as an existence claim (guardrail~G2).
  \item \textbf{Two go/no-go refinements that earned their place}
    (\cref{sec:category}, \cref{sec:parity-defect}). A categorical reading shows
    operational parity is a retraction whose round-trip defect is invisible to
    both the mutual information and the Bayes risk yet real
    (\cref{thm:retraction}); a transport/information-geometry reading metrizes
    that defect and proves it dominates the swap gap of any fixed bounded
    predictor (\cref{thm:defect-risk-gap}). Each survives only because it proves
    something the core two sections cannot (charter §2).
  \item \textbf{Retrodiction without overfitting, and falsifiable predictions}
    (\cref{sec:anchoring}, \cref{sec:predictions}). We map each of the three
    frozen Lumen cells to the regime it occupies --- T1 to the information limit,
    T3 to sub-frontier saturation, T2 to a permitted-but-unidentified bounded
    effect --- claiming only consistency, never confirmation; and we discharge
    the standing charge that a math-dressed null is unfalsifiable by issuing five
    sign-aware, falsifier-bearing predictions on the Lumen knobs
    (\cref{tab:pred-summary}).
\end{enumerate}

\subsection{Relationship to the empirical preprint, and what this paper is not}
\label{subsec:intro-companion}

\texttt{preprint-v1} is the frozen empirical record; this is its theoretical
companion, and the two are kept deliberately at arm's length. We import three
numbers from the preprint and produce none: the T1 null
($r_{\mathrm{rb}}=+0.10$, $p=0.36$), the T3 ceiling ($r_{\mathrm{rb}}=+0.43$ with
$93.7\%$ of scores saturated, effective $\nsample\approx5$), and the
Holm-absorbed T2 positive ($r_{\mathrm{rb}}=+0.58$, raw $p=0.023$, non-rejection
over the nine-cell family). Every appearance of these in \cref{sec:anchoring} is
read off the record; none is a new measurement. The theory's own contribution is
the structure around them --- the theorems, and the predictions that put the
theory at risk of being wrong.

Three disclaimers bound the whole paper and are stated once. First, the
irrelevance theorem holds at two \emph{idealizations} neither met by a frozen
LLM --- Bayes-optimality and exact parity --- and every downstream invocation
discloses which it relies on (guardrail~G3, the \texttt{\textbackslash
idealization} boxes throughout). Second, the bounded-learner objects
($\nsample$, $\hypclass$, the ERM $\widehat f_{\rep}$) are read onto a single
fixed model by analogy, so the retrodiction asserts bounds and regimes, never
point values. Third, the two go/no-go sections are provisional by charter and
carry their decision in-source; they are kept only because each earned a theorem.
What this paper is not is a claim that format does not matter, that any effect is
zero, or that the preprint should have found a positive: it is the statement that
\emph{no information-channel effect is expected}, that any real effect is a
bounded-computation phenomenon of capped size, and that the theory names where to
go looking.

% Section 2 — Preliminaries (charter §4.2; issue T1-1, WOV-262)
\section{Preliminaries: a measure-theoretic setup}
\label{sec:prelim}

% Area role (charter §2): rigorous setup so that the later DPI / sufficiency
% arguments (T2-2) are exact rather than heuristic. Everything downstream --
% information parity (§4), the bounded-learner gap (§5), the parity defect
% (§6) -- is phrased against the objects fixed here. Symbols live in
% notation.tex; this section only instantiates them.

This section fixes the measure-theoretic objects the rest of the paper
quantifies over: the data-generating measure, a \emph{representation} as a
measurable map, a \emph{predictor} as a Markov kernel, and the loss/risk
functionals. The aim is modest but load-bearing: once representations and
predictors are measurable objects on a common probability space, the
data-processing and sufficiency arguments of \cref{sec:parity} become exact
statements about conditional expectations rather than informal claims about
``information.'' We work throughout on standard Borel spaces, which is what
licences the disintegration machinery in \cref{subsec:disint}; the assumption
is harmless for our objects (finite token alphabets, countable program
syntax, and their Polish completions). We follow the standard treatments of
\citet{folland1999real}, \citet{kallenberg2002foundations}, and
\citet{cinlar2011probability}.

\subsection{The data-generating measure}
\label{subsec:dgm}

Fix a probability space $(\Omega, \mathcal{F}, \Prob)$ on which all random
elements are defined. The \emph{input space} $\Xspace$ (a program, or a
prompt-rendered code artifact) and the \emph{target space} $\Yspace$ (the
quantity to be predicted: a label, a program property, a continuation) are
each equipped with a $\sigma$-algebra; we take both to be standard Borel
spaces and write $\Bor{\Xspace}$ and $\Bor{\Yspace}$ for their Borel
$\sigma$-algebras.

\begin{assumption}[Standard Borel data model]
\label{ass:sbs}
$\Xspace$ and $\Yspace$ are standard Borel spaces. The data are a random
element $(X, Y) : \Omega \to \Xspace \times \Yspace$, measurable for
$\Bor{\Xspace} \otimes \Bor{\Yspace}$, with joint law
$\jlaw \defeq \law{(X,Y)} = \pushf{(X,Y)}{\Prob}$ on the product space. We
write $\Prob_X = \pushf{X}{\Prob}$ for the input marginal.
\end{assumption}

Standardness (every uncountable such space being Borel-isomorphic to
$[0,1]$; \citealp[Thm.~A1.2]{kallenberg2002foundations}) is used only to
guarantee that the conditional laws of \cref{subsec:disint} exist as genuine
Markov kernels. No assumption is placed on $\jlaw$ beyond measurability: in
particular the empirical regime of \texttt{preprint-v1} corresponds to a
fixed but unknown $\jlaw$, and the theory makes no claim about its form.

\subsection{Representations as measurable maps}
\label{subsec:rep}

A representation is the formatting/encoding stage that turns the raw input
into the object a predictor actually consumes (the C-ladder rungs C2--C4 and
the parity control C1$^{+}$ of \texttt{preprint-v1} are concrete instances;
their categorical relationships are deferred to \cref{sec:catrep}).

\begin{definition}[Representation]
\label{def:rep}
A \emph{representation} is a measurable map $\rep : (\Xspace, \Bor{\Xspace})
\to (\Zspace_{\rep}, \Bor{\Zspace_{\rep}})$ into a standard Borel
\emph{codomain} $\Zspace_{\rep}$. Its \emph{induced sub-$\sigma$-algebra} is
\[
  \repsig \defeq \rep^{-1}\!\big(\Bor{\Zspace_{\rep}}\big)
            \;\subseteq\; \Bor{\Xspace},
\]
and its \emph{law} is the pushforward $\pushf{\rep}{\Prob_X}$ on
$\Zspace_{\rep}$. The identity map $\mathrm{id}:\Xspace\to\Xspace$ is the
maximal (lossless) representation, with $\sigalg{\mathrm{id}} = \Bor{\Xspace}$.
\end{definition}

Two facts about \cref{def:rep} are used repeatedly. First, $\repsig$ records
exactly the information $\rep$ retains: a further map $g$ factors through
$\rep$ (i.e.\ $g = h \circ \rep$ for some measurable $h$) iff $g$ is
$\repsig$-measurable (the Doob--Dynkin lemma;
\citealp[Lem.~1.13]{kallenberg2002foundations}). Second, refinement of
representations is exactly inclusion of induced $\sigma$-algebras: if $\rep'$
factors through $\rep$ then $\sigalg{\rep'} \subseteq \repsig$. This ordering
is the backbone of the data-processing inequality in \cref{sec:parity}: a
representation can only \emph{coarsen} $\Bor{\Xspace}$, never enrich it.

\begin{remark}[Lossless vs.\ lossy]
\label{rem:lossy}
A representation is \emph{lossless for $Y$} when $\repsig$ retains all of
$X$'s predictive content for $Y$ -- formally, when $\rep(X)$ is a sufficient
statistic for $Y$ (\cref{sec:parity}, made precise in T2-1). Otherwise it is
\emph{lossy}, and the rate--distortion corollary of \cref{sec:parity} bounds
the resulting risk inflation. We do not yet need the definition of
sufficiency; we only need that $\repsig$ is the object it will be stated
against.
\end{remark}

\subsection{Predictors as Markov kernels}
\label{subsec:pred}

A frozen LLM is not a deterministic function of its input: at fixed weights
it induces, through decoding, a conditional distribution over outputs. The
right abstraction is therefore a Markov kernel, not a measurable map.

\begin{definition}[Markov kernel]
\label{def:kernel}
A \emph{Markov kernel} from $(\Zspace,\Bor{\Zspace})$ to
$(\Aspace,\Bor{\Aspace})$ is a map $\kernel{K} : \Zspace \times \Bor{\Aspace}
\to [0,1]$ such that (i) $z \mapsto \kernel{K}(z, B)$ is
$\Bor{\Zspace}$-measurable for each fixed $B \in \Bor{\Aspace}$, and (ii)
$B \mapsto \kernel{K}(z, B)$ is a probability measure on $\Bor{\Aspace}$ for
each fixed $z \in \Zspace$.
\end{definition}

\begin{definition}[Predictor; frozen LLM]
\label{def:predictor}
A \emph{predictor on the representation $\rep$} is a Markov kernel
$\kernel{K} : \Zspace_{\rep} \times \Bor{\Aspace} \to [0,1]$ from the
representation codomain to an \emph{action space} $\Aspace$ (the prediction
space; $\Aspace = \Yspace$ in the simplest case). A \emph{frozen} LLM is such
a kernel held fixed: the formatting choice is absorbed into $\rep$, and
$\kernel{K}$ is the fixed map weights-plus-decoder induces from formatted
input to output distribution.
\end{definition}

A predictor on $\rep$ acts on the raw input by pre-composition with the
representation, yielding a kernel from $\Xspace$ to $\Aspace$:
\[
  (\kcomp{\kernel{K}}{\rep})(x, B) \;\defeq\; \kernel{K}\!\big(\rep(x), B\big),
  \qquad x \in \Xspace,\; B \in \Bor{\Aspace}.
\]
Because $\rep$ is measurable and $\kernel{K}$ is a kernel,
$\kcomp{\kernel{K}}{\rep}$ is again a Markov kernel
(\citealp[Ch.~2]{cinlar2011probability}); crucially it is measurable with
respect to $\repsig$, so a predictor built on $\rep$ can never depend on input
information that $\rep$ discards. This is the measure-theoretic content of
``the predictor only sees the representation,'' and it is what the
data-processing inequality will exploit.

\subsection{Loss, risk, and Bayes risk}
\label{subsec:risk}

\begin{definition}[Loss and risk]
\label{def:risk}
A \emph{loss} is a measurable map $\loss : \Yspace \times \Aspace \to
[0,\infty]$. The \emph{risk} of a predictor $\kernel{K}$ on representation
$\rep$ is the expected loss under the data law and the predictor's own
randomness,
\[
  \risk(\kernel{K}; \rep)
  \;\defeq\;
  \Expect_{(X,Y)\sim\jlaw}\;
  \int_{\Aspace} \loss\big(Y, a\big)\,
       \kernel{K}\!\big(\rep(X), \mathrm{d}a\big),
\]
the inner integral being a Lebesgue integral against the probability measure
$\kernel{K}(\rep(X),\cdot)$.
\end{definition}

\begin{definition}[Bayes risk relative to a representation]
\label{def:bayesrisk}
The \emph{Bayes risk achievable from $\rep$} is the infimum of the risk over
all predictors that act through $\rep$,
\[
  \bayesrisk(\rep)
  \;\defeq\;
  \inf_{\kernel{K}}\, \risk(\kernel{K}; \rep),
\]
the infimum ranging over all Markov kernels
$\kernel{K} : \Zspace_{\rep} \times \Bor{\Aspace} \to [0,1]$. Equivalently
(\cref{subsec:disint}) it is the risk of the best $\repsig$-measurable
decision rule. The unconstrained Bayes risk is $\bayesrisk(\mathrm{id})$.
\end{definition}

Two remarks fix the reading of \cref{def:bayesrisk}. (a)~$\bayesrisk(\rep)$ is
the \emph{information limit} of representation $\rep$: the best any predictor
could do given only what $\rep$ retains, with capacity set aside. The whole
negative half of the paper's thesis is the statement that, under parity,
$\bayesrisk(\rep)$ does not move with $\rep$ (\cref{sec:parity}); the whole
positive half is that real, capacity-bounded predictors do not attain
$\bayesrisk(\rep)$, and that this gap -- not the information limit -- is where
format can act (\cref{sec:bounded}). (b)~Because pre-composition with $\rep$
only shrinks the available $\sigma$-algebra ($\repsig \subseteq
\Bor{\Xspace}$), the map $\rep \mapsto \bayesrisk(\rep)$ is monotone under
refinement: a coarser representation has weakly larger Bayes risk. Parity is
precisely the boundary case where the inequality is tight.

\subsection{Disintegration and the conditional-expectation machinery}
\label{subsec:disint}

The arguments of \cref{sec:parity} reduce to three classical facts. We state
them in the form used downstream; proofs are standard
(\citealp[Ch.~3,~6]{folland1999real}; \citealp[Ch.~5--6]{kallenberg2002foundations}).

\begin{proposition}[Radon--Nikodym]
\label{prop:rn}
Let $\nu$ be a $\sigma$-finite measure and $\mu \ll \nu$ a measure absolutely
continuous with respect to it on a common measurable space. Then there is a
$\nu$-a.e.\ unique nonnegative measurable density $\rnd{\mu}{\nu}$ with
$\mu(A) = \int_A \rnd{\mu}{\nu}\,\mathrm{d}\nu$ for all measurable $A$.
\end{proposition}

Radon--Nikodym is what gives conditional expectation its closed form.
For an integrable random variable $W$ and a sub-$\sigma$-algebra
$\mathcal{G} \subseteq \mathcal{F}$, the conditional expectation
$\condexp{W}{\mathcal{G}}$ is the $\Prob$-a.s.\ unique
$\mathcal{G}$-measurable variable satisfying $\Expect[\,\condexp{W}{\mathcal{G}}
\,\mathbf{1}_G\,] = \Expect[W\,\mathbf{1}_G]$ for all $G \in \mathcal{G}$;
existence is \cref{prop:rn} applied to $A \mapsto \Expect[W\,\mathbf{1}_A]$ on
$\mathcal{G}$. The single property we use is the \emph{tower rule}: for nested
$\sigma$-algebras $\mathcal{G}_1 \subseteq \mathcal{G}_2$,
\begin{equation}
  \label{eq:tower}
  \condexp{\condexp{W}{\mathcal{G}_2}}{\mathcal{G}_1}
  = \condexp{W}{\mathcal{G}_1}.
\end{equation}
Applied to $\mathcal{G}_1 = \repsig \subseteq \mathcal{G}_2 = \Bor{\Xspace}$,
\cref{eq:tower} is the data-processing inequality in embryo: coarsening to
$\repsig$ is a conditional-expectation projection, and projections cannot
increase the information available to a convex risk. \cref{sec:parity} turns
this into the irrelevance theorem.

\begin{proposition}[Disintegration / regular conditional probability]
\label{prop:disint}
Under \cref{ass:sbs}, for the sub-$\sigma$-algebra $\repsig$ there exists a
Markov kernel $\post{\rep} : \Zspace_{\rep} \times \Bor{\Yspace} \to [0,1]$,
the \emph{regular conditional law of $Y$ given $\rep(X)$}, such that
\[
  \post{\rep}\!\big(\rep(X), B\big)
  = \Prob\big(Y \in B \given \repsig\big)
  \qquad \Prob\text{-a.s., for every } B \in \Bor{\Yspace}.
\]
It is $\pushf{\rep}{\Prob_X}$-a.s.\ unique. The kernel $\post{\rep}$ is the
\emph{Bayes-optimal predictor} on $\rep$: for any loss $\loss$ admitting a
measurable Bayes act, plugging $\post{\rep}$ into \cref{def:risk} attains
$\bayesrisk(\rep)$ of \cref{def:bayesrisk}.
\end{proposition}

Existence in \cref{prop:disint} is exactly where standardness of $\Yspace$ is
spent (\citealp[Thm.~6.3]{kallenberg2002foundations};
\citealp[Ch.~IV]{cinlar2011probability}); on a non-standard target the
conditional law need not regularize into a kernel. \cref{prop:disint} closes
the loop opened in \cref{subsec:pred}: the abstract ``best predictor on
$\rep$'' of \cref{def:bayesrisk} is a concrete Markov kernel, the posterior
$\post{\rep}$, and comparing $\post{\rep}$ across representations under parity
is precisely the calculation \cref{sec:parity} carries out. Everything later
in the paper is an estimate of how far a frozen LLM's kernel $\kernel{K}$
sits from $\post{\rep}$, and of how that distance moves -- or, under parity,
fails to move -- with $\rep$.

% Section 3 — The category of representations (charter §4.3; issues T1-4, T2-3)
% GO/NO-GO area (charter §2): category theory survives only if it earns a theorem.
\section{The category of representations}
\label{sec:category}

\gonogo{Category theory (T1-4 $\to$ T2-3) is provisional. If the
parity-as-retraction framing produces only a restatement of the
information-theoretic content of \cref{sec:parity} and buys no new
proposition, this section is demoted to a remark and the decision is
recorded in WOV-269 (T2-3).}

% STUB (T1-4, T2-3). The C-ladder as morphisms; information parity as a
% retraction; the round-trip defect \roundtrip{\ladder{4}}{\ladderplus}.
% Area role (charter §2): why C1+ is the correct control (parity = retraction).
%
% Authoritative C-ladder (charter §5; design-constitution.tex §14.1):
%   C2 = Raw AST, C3 = Typed AST, C4 = Contract IR,
%   C1+ = annotated source projected from C4.
% Use these definitions, NOT the README's simplified table.

\noindent\emph{[Stub. Fills in T1-4 (foundations) and T2-3 (the
retraction proposition, if it earns its place).]}

% Section 4 — Information parity and the irrelevance theorem
% (charter §4.4; issues T1-2, T2-1, T2-2)
\section{Information parity and the irrelevance theorem}
\label{sec:parity}

% STUB (T1-2, T2-1, T2-2). The spine of the paper.
% Area role (charter §2): parity = sufficiency / Markov-equivalence; the
% irrelevance theorem via DPI; rate-distortion for lossy reps.
%
% Plan:
%   - T2-1: formal definitions of structured representation and information
%     parity (parity = \rep(X) \suff X w.r.t. Y, i.e. X \markov \rep(X) \markov Y
%     with no loss of sufficiency).
%   - T2-2: the at-optimum invariance theorem: under parity, \bayesrisk is
%     invariant to \rep. DPI is the engine.
%   - Corollary: rate-distortion bound \ratedist for lossy representations.

\noindent\emph{[Stub. Definitions in T2-1; irrelevance theorem in T2-2;
information-theoretic foundations in T1-2.]}

\begin{theorem}[Irrelevance at optimum --- placeholder]
\label{thm:irrelevance}
\emph{[Placeholder for T2-2.]} Under information parity
($X \markov \rep(X) \markov Y$ with $\rep(X) \suff X$), the Bayes risk is
format-invariant: $\bayesrisk(\rep) = \bayesrisk(\mathrm{id})$.
\idealization{This statement assumes Bayes-optimality (the
infinite-capacity limit) and exact parity. A real frozen LLM satisfies
neither; \cref{sec:parity-defect} quantifies the parity defect and
\cref{sec:bounded} relaxes the capacity idealization. ``Parity is
operational, not absolute'' (preprint §3.3) is carried in here as a formal
caveat.}
\end{theorem}

% Section 5 — The bounded-learner gap (charter §4.5; issues T1-3, T2-4)
\section{The bounded-learner gap}
\label{sec:bounded}

% STUB (T1-3, T2-4). The positive half of the thesis (charter §1).
% Area role (charter §2): where format can help once information is held
% constant; an UPPER BOUND on achievable representation advantage via
% capacity/sample complexity; the regime that activates it.
%
% Plan:
%   - T1-3: statistical learning theory foundations (hypothesis class
%     \hypclass, capacity \capac{\hypclass}, Rademacher \Rad{\hypclass},
%     empirical risk \emprisk).
%   - T2-4: the bounded-learner advantage upper bound: the representation
%     advantage \advantage is controlled by a capacity/complexity term, not
%     by information content. This bounds effect SIZE and is the antidote to
%     "a math-dressed null" (guardrail G2).
%
% The positive content is delivered as falsifiable predictions in §8, never
% as established findings (guardrail G2).

\noindent\emph{[Stub. Learning-theory foundations in T1-3; advantage upper
bound in T2-4.]}

\begin{proposition}[Bounded-regime advantage is capacity-controlled --- placeholder]
\label{prop:advantage}
\emph{[Placeholder for T2-4.]} The achievable representation advantage
$\advantage$ at finite capacity satisfies an upper bound in terms of
$\capac{\hypclass}$ and the sample size $\nsample$.
\idealization{Holds in the bounded regime (finite capacity / finite sample
/ bounded compute); it does not assert any effect exists, only bounds its
maximum magnitude.}
\end{proposition}

% Section 6 — Quantitative parity defect (charter §4.6; issues T1-5, T2-5)
% GO/NO-GO area (charter §2): OT/info-geometry survives only if it earns a theorem.
\section{Quantitative parity defect}
\label{sec:parity-defect}

\gonogo{Optimal transport + information geometry (T1-5 $\to$ T2-5) was
provisional. If the parity defect yielded only a restatement of the
information-theoretic content of \cref{sec:parity} and bought no new
proposition, this section was to be demoted to a remark and the decision
recorded in WOV-271 (T2-5). T1-5 supplied the foundations and \emph{defined}
the metric precisely enough to be computed against the R4 artifacts; the
explicit criterion is stated in \cref{subsec:defect-gonogo}.
\textbf{Resolved (T2-5, WOV-271): go.} The metric earns
\cref{thm:defect-risk-gap}: $\paritydefect$ upper-bounds the risk gap of a
\emph{fixed} bounded predictor across the
$\ladder{4}\!\leftrightarrow\!\ladderplus$ swap --- a quantity
\cref{thm:irrelevance} and the information gap $\infogap{\cdot}$ cannot see
(both vanish under parity) --- and \cref{cor:defect-finer-order} shows it
strictly refines the mutual-information ordering on the parity class, witnessed
by the R4 profile (\cref{ex:r4-defect}). The section stands
(\cref{subsec:defect-theorem}).}

% This file holds the T1-5 (WOV-266) FOUNDATIONS only: Wasserstein distance +
% Kantorovich-Rubinstein duality; statistical manifolds, the Fisher metric,
% KL as a local divergence; the DEFINITION of the quantitative parity defect
% (W1 between induced predictive distributions + a local Fisher sensitivity);
% and a recipe for computing it in principle against the R4 round-trip data.
% The metric is USED (and earns or loses its place) in T2-5.

\cref{sec:parity} treats parity as a binary equality of mutual informations,
and guardrail~G3 flags that real controls meet it only ``operationally, not
absolutely.'' This section turns ``how far from parity'' into a number. The
mutual-information gap $\MI{Y}{X}-\MI{Y}{\rep(X)}$ already measures lost
\emph{information}; what it does not measure is the residual asymmetry between
two representations that carry the \emph{same} information about $Y$ yet drive
a bounded predictor to slightly different predictive distributions --- exactly
the $\sigalg{\ladder{4}}\ne\sigalg{\ladderplus}$ gap of
\cref{prop:abs-strict}. We build the optimal-transport and
information-geometry tools that quantify that residual, define the parity
defect $\paritydefect$, and specify how it is computed against the R4
round-trip artifacts. We follow \citet{villani2009} and \citet{peyrecuturi2019}
for optimal transport and \citet{amarinagaoka2000} for information geometry.

\subsection{Wasserstein distance and Kantorovich--Rubinstein duality}
\label{subsec:wasserstein}

\begin{definition}[Wasserstein distance]
\label{def:wasserstein}
Let $(\Aspace,d)$ be a Polish metric space and $\mu,\nu$ probability measures
on it with finite first moment. The \emph{$1$-Wasserstein distance} is
\[
  \Wass{1}(\mu,\nu)
  \;\defeq\;
  \inf_{\gamma\in\Pi(\mu,\nu)}\;
  \int_{\Aspace\times\Aspace} d(a,a')\,\mathrm{d}\gamma(a,a'),
\]
the infimum over couplings $\gamma$ with marginals $\mu$ and $\nu$. It is the
minimal average cost of transporting $\mu$ onto $\nu$.
\end{definition}

\begin{proposition}[Kantorovich--Rubinstein duality]
\label{prop:kr-duality}
With $\mu,\nu$ as above,
\begin{equation}
  \label{eq:kr}
  \Wass{1}(\mu,\nu)
  \;=\;
  \sup_{\,\|g\|_{\mathrm{Lip}}\le 1}\;
  \Big(\,\Expect_{\mu}[g] - \Expect_{\nu}[g]\,\Big),
\end{equation}
the supremum over $1$-Lipschitz test functions $g:\Aspace\to\mathbb{R}$.
\end{proposition}

\begin{proof}
This is the dual form of the Kantorovich problem for the cost $d$; see
\citet[Thm.~5.10 and Particular Case~5.16]{villani2009}.
\end{proof}

The duality \eqref{eq:kr} is what makes the defect \emph{estimable}: it
rewrites a transport cost over couplings as an expectation gap over Lipschitz
witnesses, which can be evaluated from samples of the two predictive
distributions without solving the primal coupling problem.

\subsection{Statistical manifolds and the Fisher metric}
\label{subsec:fisher}

Where the Wasserstein distance measures the global cost of moving one
predictive distribution to another, the Fisher metric measures the
\emph{local} sensitivity of a predictor to small representation
perturbations, which is the regime relevant when $\ladderplus$ and
$\ladder{4}$ are close.

\begin{definition}[Statistical manifold, Fisher metric]
\label{def:fisher}
A parametric family $\{p_\theta:\theta\in\Theta\}$ of densities, smooth in
$\theta$ over an open $\Theta\subseteq\mathbb{R}^k$, is a \emph{statistical
manifold}. Its \emph{Fisher information metric} at $\theta$ is the
$k\times k$ matrix
\[
  \Fisher(\theta)_{ij}
  \;\defeq\;
  \Expect_{p_\theta}\!\Big[
    \partial_{\theta_i}\log p_\theta \;
    \partial_{\theta_j}\log p_\theta
  \Big],
\]
the covariance of the score. It is the unique (up to scale) Riemannian metric
invariant under sufficient statistics \citep[Ch.~2--3]{amarinagaoka2000}.
\end{definition}

\begin{proposition}[KL is the Fisher metric to second order]
\label{prop:kl-fisher}
For $\theta'=\theta+\mathrm{d}\theta$,
\[
  \KL{p_\theta}{p_{\theta'}}
  \;=\;
  \tfrac12\,\mathrm{d}\theta^\top \Fisher(\theta)\,\mathrm{d}\theta
  \;+\; o(\|\mathrm{d}\theta\|^2).
\]
\end{proposition}

\begin{proof}
Taylor-expand $\theta'\mapsto\KL{p_\theta}{p_{\theta'}}$ at $\theta'=\theta$:
the value and gradient vanish, and the Hessian is $\Fisher(\theta)$
\citep[Ch.~3]{amarinagaoka2000}.
\end{proof}

\cref{prop:kl-fisher} ties this section back to \cref{sec:parity}: the KL
divergence that defines mutual information (\cref{def:mi}) is, locally, the
Fisher quadratic form. So a Fisher-sensitivity term and a relative-entropy
term are the same object at two scales --- which is why the defect below pairs
a global Wasserstein term with a local Fisher term rather than choosing
between them.

\subsection{The quantitative parity defect}
\label{subsec:defect-def}

Let $\preddist{c}$ denote the \emph{induced predictive distribution} under
condition $c$: the law on the action space $\Aspace$ of the bounded
predictor's output $\kernel{K}(\rep_c(X),\cdot)$ when $X\sim\Prob_X$, for
$c\in\{\ladder{4},\ladderplus\}$. Under \emph{exact} parity these two laws
coincide for the Bayes predictor; for a bounded predictor they may differ even
when $\MI{Y}{\ladder{4}(X)}=\MI{Y}{\ladderplus(X)}$.

\begin{definition}[Quantitative parity defect]
\label{def:parity-defect}
Fix a bounded predictor $\kernel{K}$ and a perturbation scale. The
\emph{parity defect} between $\ladder{4}$ and $\ladderplus$ is
\begin{equation}
  \label{eq:parity-defect}
  \paritydefect
  \;\defeq\;
  \underbrace{\Wass{1}\!\big(\preddist{\ladder{4}},\,\preddist{\ladderplus}\big)}_{\text{global: predictive-law transport}}
  \;+\;
  \lambda\,
  \underbrace{\sqrt{\,\overline{v}^\top \Fisher\, \overline{v}\,}}_{\text{local: Fisher sensitivity}},
\end{equation}
where $\overline{v}$ is the mean representation perturbation induced by the
round trip $\roundtrip{\ladder{4}}{\ladderplus}$ (\cref{def:cat-parity}),
identified by \cref{cor:defect-subadditive} with the complement of the
idempotent splitting of \cref{thm:retraction},
$\Fisher$ is the Fisher metric of the predictor's output family at the
operating point, and $\lambda\ge0$ trades off the two scales. Exact parity is
$\paritydefect=0$; $\paritydefect>0$ certifies operational-but-not-absolute
parity and bounds, via \cref{prop:kr-duality}, the largest difference in
expected Lipschitz functionals of the two predictive distributions.
\end{definition}

The two terms are complementary by \cref{prop:kl-fisher}: the Wasserstein term
captures the part of the asymmetry that survives at the global scale of the
output distribution, and the Fisher term captures the local sensitivity of the
predictor to the residual structural perturbation that the round trip fails to
undo. \cref{eq:parity-defect} is a definition, not yet a theorem; T2-5 must
show $\paritydefect$ either bounds a risk gap or constrains the predictive
asymmetry in a way the mutual-information gap of \cref{sec:parity} does not
(\cref{subsec:defect-gonogo}).

\subsection{Computing the defect against the R4 round-trip data}
\label{subsec:defect-r4}

The R4 round-trip recovery experiment (empirical paper Appendix~D) re-parses
$\ladderplus$ back to a Contract IR $\ladder{4}'$ and reports how much of
$\ladder{4}$'s structure is recovered: contracts $100\%$, AST node kinds
$99.8\%$, types $96.4\%$, parent--child edges $98.7\%$, byte-exact AST
$0/30$. This profile is, in the categorical reading of
\cref{subsec:retraction-parity}, the empirical signature of the idempotent
$e\circ p$: near-identity on semantic structure, far from identity on surface
bytes. The defect \eqref{eq:parity-defect} is computed against it as follows.

\begin{enumerate}
  \item \textbf{Perturbation $\overline{v}$ (Fisher term).} Each
    non-recovered structural element of the round trip is a coordinate of the
    representation perturbation; the recovery rates above estimate the
    per-coordinate failure probabilities, and $\overline{v}$ is their
    aggregate over the $30$ functions. Byte-level coordinates (recovery $0/30$)
    dominate $\overline{v}$ in magnitude but are predictively inert, which is
    why the Fisher term --- weighting perturbations by their effect on the
    \emph{output} family --- and not the raw byte distance is the right local
    measure.
  \item \textbf{Predictive laws (Wasserstein term).} Sample the bounded
    predictor on the C4 inputs and on their re-parsed $\ladder{4}'$
    counterparts to obtain empirical $\preddist{\ladder{4}}$ and
    $\preddist{\ladderplus}$, then estimate
    $\Wass{1}(\preddist{\ladder{4}},\preddist{\ladderplus})$ either by the
    primal optimal-transport solver of \citet[§2--3]{peyrecuturi2019} or by
    maximizing the Kantorovich--Rubinstein dual \eqref{eq:kr} over a Lipschitz
    witness class.
  \item \textbf{Trade-off $\lambda$.} Fixed by matching the two terms' units on
    a held-out calibration slice; the procedure is specified with the T2-5
    computation, not chosen post hoc.
\end{enumerate}

\begin{remark}[Status of the computation, G3/G6]
\label{rem:defect-status}
The recipe is ``computable in principle'': it specifies estimable quantities
and named estimators, but it is \emph{not} executed in this paper, and no
numerical value of $\paritydefect$ is claimed. The R4 rates are used as the
inputs the metric \emph{would} consume, illustrating that the defect is
grounded in measured artifacts rather than postulated. Any statement that the
computed defect is small (and hence that the empirical control was close to
exact parity) is, until T2-5 runs it, a conjecture and is labeled as such
(guardrail~G6).
\end{remark}

\subsection{Go/no-go criterion handed to T2-5}
\label{subsec:defect-gonogo}

\begin{quote}
\textbf{Criterion (depth of OT/info-geometry).} Section~\ref{sec:parity-defect}
survives as a section iff T2-5 establishes that $\paritydefect$
\eqref{eq:parity-defect} (i) upper-bounds the bounded predictor's risk gap
between $\ladder{4}$ and $\ladderplus$, or (ii) provides a strictly finer
ordering of representations than the mutual-information gap of
\cref{sec:parity} (two reps with equal information gap but different defect,
witnessed against the R4 profile). If T2-5 produces only a reformulation of
the information gap, the section is demoted to a remark recording that the
parity defect coincides with the information gap in this setting, and the
decision is logged in WOV-271.
\end{quote}

\subsection{What the metric earns: the defect bounds the bounded-predictor risk gap (T2-5)}
\label{subsec:defect-theorem}

This subsection discharges the go/no-go of \cref{subsec:defect-gonogo}, and the
verdict is \textbf{go}. The criterion asks for \emph{either} prong; the defect
delivers both. We show (i) $\paritydefect$ upper-bounds the risk gap of a fixed
bounded predictor across the $\ladder{4}\leftrightarrow\ladderplus$ swap
(\cref{thm:defect-risk-gap}, with the Fisher term supplying the local regime in
\cref{prop:defect-fisher-local}), and (ii) it strictly refines the
mutual-information ordering on the parity class
(\cref{cor:defect-finer-order}), a separation realized by the R4 control pair
(\cref{ex:r4-defect}). Full proofs are in \cref{app:parity-defect}; the decision
is recorded in WOV-271.

The object both prongs concern is the \emph{shared reader}: a single bounded
predictor consumes either representation. Formally we keep exactly the setup
\cref{def:parity-defect} already commits to --- one kernel $\kernel{K}$ evaluated
at $\rep_c(X)$ for $c\in\{\ladder{4},\ladderplus\}$ --- which models the frozen
LLM of \texttt{preprint-v1} reading $\ladder{4}$ or $\ladderplus$ through their
common token-sequence codomain. Its risk on condition $c$ is
$\risk(\kernel{K};c)=\Expect_{(X,Y)}\!\int_{\Aspace}\loss(Y,a)\,\kernel{K}(\rep_c(X),\mathrm{d}a)$
(\cref{def:risk}), and the quantity at issue is the
\emph{bounded-predictor representation gap}
\begin{equation}
  \label{eq:bp-gap}
  \advantage_{\kernel{K}}(\ladder{4},\ladderplus)
  \;\defeq\;
  \risk(\kernel{K};\ladder{4})-\risk(\kernel{K};\ladderplus).
\end{equation}
This is \emph{not} the advantage $\advantage(\rep,\rep')$ of \cref{def:advantage}:
that compares two \emph{independently optimized} ERMs and is capacity-bounded
(\cref{thm:bounded-gap}); \eqref{eq:bp-gap} compares \emph{one fixed} predictor
across the two formats and is, as we now show, transport-bounded. The two are
different operational quantities, which is exactly why \cref{sec:parity-defect}
is not subsumed by \cref{sec:bounded}.

\begin{definition}[Input-coupled predictive transport]
\label{def:coupled-transport}
For the shared predictor $\kernel{K}$ define
\begin{equation}
  \label{eq:coupled-transport}
  \cpredtransport
  \;\defeq\;
  \Expect_X\!\Big[\,
    \Wass{1}\big(\kernel{K}(\ladder{4}(X),\cdot),\,\kernel{K}(\ladderplus(X),\cdot)\big)
  \,\Big],
\end{equation}
the input-coupled form of the global term of \eqref{eq:parity-defect}. By joint
convexity of $\Wass{1}$ \citep[Thm.~4.8]{villani2009} and
$\preddist{c}=\Expect_X[\kernel{K}(\rep_c(X),\cdot)]$,
\begin{equation}
  \label{eq:marginal-le-coupled}
  \Wass{1}\big(\preddist{\ladder{4}},\preddist{\ladderplus}\big)\;\le\;\cpredtransport,
\end{equation}
with equality when the input pairing already realizes an optimal coupling of the
marginals (e.g.\ a swap-insensitive predictor). The marginal term of
\cref{def:parity-defect} is thus the relaxation; $\cpredtransport$ is the
risk-controlling form, and it inherits the rung-wise subadditivity of
\cref{cor:defect-subadditive} ($\Expect_X\Wass{1}$ is again a metric on the
predictive laws).
\end{definition}

\begin{theorem}[The defect dominates the bounded-predictor risk gap]
\label{thm:defect-risk-gap}
Let $\kernel{K}$ be a fixed bounded predictor read on both representations and
suppose the per-instance loss $a\mapsto\loss(y,a)$ is $L$-Lipschitz on
$(\Aspace,d)$ uniformly in $y$. Then
\begin{equation}
  \label{eq:defect-risk-gap}
  \big|\advantage_{\kernel{K}}(\ladder{4},\ladderplus)\big|
  \;\le\;
  L\,\cpredtransport,
\end{equation}
the global (Wasserstein) term of the parity defect \eqref{eq:parity-defect}.
Under operational parity $\ladder{4}\parityeq\ladderplus$ the information gap
contributes nothing --- $\infogap{\ladder{4}}=\infogap{\ladderplus}=0$
(\cref{prop:suff-iff-mi}), so the middle term of the excess-risk decomposition
\eqref{eq:excess} cancels between the two conditions --- and the entire surviving
gap of the shared predictor is controlled by the transport defect, with no
mutual-information term.
\idealization{This relaxes idealization~(i) of \cref{rem:slogan}
(Bayes-optimality): $\kernel{K}$ is an arbitrary fixed bounded predictor, not the
posterior $\post{\rep}$. It engages the failure of idealization~(ii) in its sharp
form --- the bound is informative precisely when
$\sigalg{\ladder{4}}\ne\sigalg{\ladderplus}$ (\cref{prop:abs-strict}) makes
$\kernel{K}$ respond to the swap; under absolute parity both sides are $0$. The
bound limits the \emph{size} of the swap gap in either direction and asserts no
effect exists (guardrail~G2); the Lipschitz-loss hypothesis is discharged in the
near regime by \cref{prop:defect-fisher-local} for merely bounded losses.}
\end{theorem}

\begin{proof}[Proof idea]
For fixed $(x,y)$, $\loss(y,\cdot)/L$ is $1$-Lipschitz, so
Kantorovich--Rubinstein duality \eqref{eq:kr} bounds the per-instance gap of
expected loss by $L\,\Wass{1}(\kernel{K}(\ladder{4}(x),\cdot),
\kernel{K}(\ladderplus(x),\cdot))$, uniformly in $y$. Take $\Expect_{(X,Y)}$ and
pull the absolute value inside. The information-gap cancellation is
\eqref{eq:excess} read for the shared kernel under \cref{prop:suff-iff-mi}. Full
proof in \cref{app:parity-defect}.
\end{proof}

The Lipschitz hypothesis is the right one when the two predictive laws are far
apart, but it says nothing when the loss is merely bounded (the regime of
\cref{cor:ratedist-gap}(b)). There the \emph{local} Fisher term takes over,
exactly as \cref{prop:kl-fisher} anticipated.

\begin{proposition}[Local Fisher control of the swap gap]
\label{prop:defect-fisher-local}
Suppose instead $\loss$ is bounded, $\|\loss\|_\infty<\infty$, and the predictor's
output family is smooth in the representation perturbation, the round trip
$\roundtrip{\ladder{4}}{\ladderplus}$ inducing the mean perturbation
$\overline{v}$ of \cref{def:parity-defect}. Then
\begin{equation}
  \label{eq:defect-fisher-local}
  \big|\advantage_{\kernel{K}}(\ladder{4},\ladderplus)\big|
  \;\le\;
  \|\loss\|_\infty\,\sqrt{\tfrac14\,\overline{v}^\top \Fisher\, \overline{v}}
  \;+\; o(\|\overline{v}\|),
\end{equation}
so to leading order the swap gap is governed by the Fisher term of
\eqref{eq:parity-defect}.
\idealization{Like \cref{thm:defect-risk-gap} this is a bounded-predictor
statement (idealization~(i) relaxed). It is a local (small-$\overline{v}$)
bound: the $o(\|\overline{v}\|)$ remainder is the Taylor error of
\cref{prop:kl-fisher}, and the inequality is informative only when the round-trip
perturbation is small --- the near-parity regime in which the global bound
\eqref{eq:defect-risk-gap} may be loose. Sign-agnostic (G2).}
\end{proposition}

\begin{proof}[Proof idea]
Per instance, the bounded-loss/total-variation step of
\cref{cor:ratedist-gap}(b) gives
$|\!\int\loss(y,\cdot)\,\mathrm{d}(p-p')|\le\|\loss\|_\infty\|p-p'\|_{\mathrm{TV}}$;
Pinsker bounds $\|p-p'\|_{\mathrm{TV}}\le\sqrt{\tfrac12\KL{p}{p'}}$ and
\cref{prop:kl-fisher} expands
$\KL{p}{p'}=\tfrac12\overline{v}^\top\Fisher\overline{v}+o(\|\overline{v}\|^2)$ at
the aggregate increment $\overline{v}$ of \cref{def:parity-defect}. The bound is
then $x$-free, so $\Expect_{(X,Y)}$ leaves it unchanged. Full proof in
\cref{app:parity-defect}.
\end{proof}

\begin{corollary}[The two terms are complementary, not redundant]
\label{cor:defect-two-term}
With $\lambda$ fixed by the unit-matching calibration of \cref{subsec:defect-r4},
there is a constant $c=c(L,\|\loss\|_\infty,\lambda)$ with
\[
  \big|\advantage_{\kernel{K}}(\ladder{4},\ladderplus)\big|
  \;\le\; c\,\paritydefect \;+\; o(\|\overline{v}\|),
\]
the global term \eqref{eq:defect-risk-gap} dominating the far regime (Lipschitz
loss, predictive laws far apart) and the Fisher term
\eqref{eq:defect-fisher-local} the near regime (bounded loss, small
perturbation). Each term bounds the swap gap exactly where the other is loose or
unavailable, which is why \eqref{eq:parity-defect} carries both rather than
choosing between them.
\end{corollary}

\begin{proof}
Take the smaller of the two right-hand sides of \eqref{eq:defect-risk-gap} and
\eqref{eq:defect-fisher-local} and bound it by their sum, absorbing
$L,\|\loss\|_\infty,\lambda$ into $c$.
\end{proof}

This settles prong~(i). Prong~(ii) is now immediate and is what makes the defect
strictly more than a re-encoding of \cref{sec:parity}.

\begin{corollary}[The defect strictly refines the information-gap ordering]
\label{cor:defect-finer-order}
On the class of representations in operational parity with $\ladder{4}$, the
information gap $\infogap{\cdot}$ is constant ($\equiv 0$,
\cref{prop:suff-iff-mi}) and so, by \cref{thm:irrelevance}, is the Bayes risk:
neither orders the class. The parity defect does. It is generically positive
there --- $\paritydefect>0$ whenever $\sigalg{\ladder{4}}\ne\sigalg{\ladderplus}$
(\cref{prop:abs-strict}) and $\kernel{K}$ is swap-sensitive, the case
\cref{thm:retraction}(iii) places in the common kernel of $\infogap{\cdot}$ and
$\bayesrisk(\cdot)$ --- and by \cref{thm:defect-risk-gap} it ranks parity pairs by
the bounded-predictor risk gap they admit. Hence $\paritydefect$ separates
representations the information gap and the Bayes risk cannot, a strict
refinement of both.
\end{corollary}

\begin{proof}
Constancy of $\infogap{\cdot}$ and $\bayesrisk(\cdot)$ on the parity class is
\cref{prop:suff-iff-mi} and \cref{thm:irrelevance}. Positivity of
$\paritydefect$ off absolute parity is \cref{def:parity-defect} with
\cref{prop:abs-strict}; that it bounds (hence orders by) the swap gap is
\cref{thm:defect-risk-gap}. The control pair
$(\ladder{4},\ladderplus)$ has $\infogap{\ladder{4}}=\infogap{\ladderplus}=0$ by
design yet $\paritydefect>0$ by \cref{ex:r4-defect}, an explicit witness.
\end{proof}

\begin{example}[The defect on the R4 round trip: an illustrative computation]
\label{ex:r4-defect}
We instantiate \eqref{eq:parity-defect} on the R4 profile by the recipe of
\cref{subsec:defect-r4}. Each non-recovered structural element is a coordinate of
$\overline{v}$ with magnitude its round-trip failure rate ($1-$recovery);
\cref{tab:r4-defect} collects them. The Fisher term weights each coordinate by
its effect on the predictor's output: the byte-exact coordinate, recovered
$0/30$ and hence maximal in raw magnitude, is \emph{predictively inert} and
carries Fisher weight $\approx 0$, while the semantic coordinates carry unit
weight (an illustrative diagonal $\Fisher$). With the action diameter normalized
to $1$ and $L=\lambda=1$, the global term is bounded by the predictively-relevant
failure mass and the local term by the Fisher quadratic form:
\[
  \cpredtransport \;\lesssim\; 0.051,
  \qquad
  \sqrt{\overline{v}^\top\Fisher\,\overline{v}} \;=\; 0.038,
  \qquad
  \paritydefect \;\approx\; 0.089 .
\]
The contrast is the point. A naive defect that used raw structural distance ---
unit weight on \emph{every} coordinate, byte coordinate included --- would read
$\|\overline{v}\|\approx 1.00$, larger by a factor of $\approx 26$ and dominated
entirely by the byte-recovery catastrophe. The Fisher weighting is exactly what
collapses that $1.00$ to $0.04$, certifying the control as \emph{near} parity
despite $0/30$ byte recovery. This is work neither the raw byte distance (which
sees a near-maximal perturbation) nor the information gap (which is identically
$0$) can do, and it is the quantitative form of the qualitative statement of
\cref{subsec:retraction-parity} that the round-trip idempotent is near-identity
on semantic structure and far from it on surface bytes.
\end{example}

\begin{table}[t]
\centering
\small
\begin{tabular}{@{}lrrl@{}}
\toprule
Round-trip coordinate & Recovery & $\overline{v}_i$ ($=1-$rec.) & Predictive role \\
\midrule
Contracts            & $100\%$  & $0.000$ & relevant (weight $1$) \\
AST node kinds       & $99.8\%$ & $0.002$ & relevant (weight $1$) \\
Parent--child edges  & $98.7\%$ & $0.013$ & relevant (weight $1$) \\
Types                & $96.4\%$ & $0.036$ & relevant (weight $1$) \\
Byte-exact AST       & $0/30$   & $1.000$ & \emph{inert} (weight $\approx 0$) \\
\midrule
\multicolumn{2}{@{}l}{Global term $\cpredtransport\lesssim\sum_{\text{rel.}}\overline{v}_i$} & $0.051$ & union-bound proxy \\
\multicolumn{2}{@{}l}{Fisher term $\sqrt{\overline{v}^\top\Fisher\overline{v}}$} & $0.038$ & semantic coords only \\
\multicolumn{2}{@{}l}{Illustrative $\paritydefect$ ($L=\lambda=1$)} & $0.089$ & \eqref{eq:parity-defect} \\
\bottomrule
\end{tabular}
\caption{The parity defect computed on the R4 round-trip recovery profile
(\cref{subsec:defect-r4}, \cref{ex:r4-defect}). Numbers are \emph{illustrative}:
action diameter normalized to $1$, $L=\lambda=1$, and a diagonal Fisher metric
with unit weight on the predictively-relevant coordinates and zero on the inert
byte coordinate. They are not the output of running the frozen predictor
(guardrail~G6; \cref{rem:defect-status}). The naive raw-distance defect
$\|\overline{v}\|\approx 1.00$ is $\approx 26\times$ larger and byte-dominated;
the Fisher weighting is what makes the defect small.}
\label{tab:r4-defect}
\end{table}

\begin{remark}[Why this passes the criterion; scope of the go (G4/G6)]
\label{rem:defect-go-scope}
\cref{thm:defect-risk-gap} and \cref{cor:defect-finer-order} meet the criterion of
\cref{subsec:defect-gonogo} on \emph{both} prongs (it required either). The defect
is not a reformulation of the information gap because the quantity it controls ---
the swap risk gap \eqref{eq:bp-gap} of a fixed bounded predictor --- is provably
in the common kernel of $\infogap{\cdot}$ and $\bayesrisk(\cdot)$
(\cref{thm:retraction}(iii)): both are $0$ across the parity pair while the defect
is positive (\cref{ex:r4-defect}). So a proposition that bounds and orders by the
swap gap cannot be a restatement of either. The two terms are shown
non-redundant by \cref{cor:defect-two-term}, and the R4 example shows the Fisher
term in particular does work the raw byte distance cannot. \textbf{Scope (G6).}
The numeric defect of \cref{ex:r4-defect} is illustrative and computed-in-principle
(\cref{rem:defect-status}): it consumes the measured R4 rates but normalizes the
action geometry, fixes $L,\lambda$, and posits a diagonal Fisher metric rather
than estimating one from the frozen predictor's outputs. The go rests on the
propositions \cref{thm:defect-risk-gap}, \cref{prop:defect-fisher-local},
\cref{cor:defect-finer-order}; the R4 example anchors them in artifacts and
exhibits the prong-(ii) separation, and no claim is made that the executed value
is small until T2-5's recipe is actually run (a conjecture, so labeled).
\end{remark}

% Section 7 — Empirical anchoring
\section{Empirical anchoring}
\label{sec:anchoring}

The preceding sections are self-contained mathematics: \cref{thm:irrelevance}
closes the information channel at the optimum, \cref{thm:bounded-gap} bounds
what remains for a bounded learner, and \cref{sec:parity-defect} measures the
residual when parity is only operational. This section connects those results
to the three frozen outcomes of \texttt{preprint-v1} --- the T1 clean null, the
T3 ceiling, and the T2 moderate-but-absorbed positive --- by naming, for each,
the theoretical regime it falls in. The exercise is strictly
\emph{retrodiction}: every empirical number here is read off the frozen record,
none is produced by the theory, and no new empirical claim is built. The point
is to show that the recorded outcomes are the ones the theory
\emph{already permits}, and to be exact about where the mapping from an
idealized learner to a single frozen LLM stops being a deduction and becomes an
analogy.

Three caveats bind the whole section and are stated once here rather than
repeated. First, the frozen artifact is one fixed model queried under five
representations, not an ERM ranging over a hypothesis class with a tunable
sample size $\nsample$; the PAC/Rademacher objects of \cref{sec:bounded} are
therefore read onto it by analogy, with the model's in-context behaviour
standing in for $\widehat f_{\rep}$ and the task corpus for the draw $S$
(disclosed again at each use). Second, the control is C1${}^{+}$ at
\emph{operational} parity (\cref{def:parity}), not absolute parity, so every
invocation of \cref{thm:irrelevance} below inherits idealization~(ii) of
\cref{rem:slogan} and is corrected by the \cref{sec:parity-defect} defect.
Third, none of what follows is a test: the theory is consistent with the record,
which is necessary but not sufficient for the theory to be true.

\subsection{The T1 null as the predicted information-limit outcome}
\label{subsec:anchor-t1}

The central H1 test on T1 --- the cleanest unconfounded cell, with no rubric
ceiling --- recorded a rank-biserial effect $r_{\mathrm{rb}}=+0.10$ at raw
$p=0.36$, a clean non-rejection of C4${}\le{}$C1${}^{+}$. Under the theory this
is not a weak or disappointing positive; it is the \emph{predicted}
outcome of the information-limit regime, and that prediction is exactly what
the honesty discipline forbids us to soften.

The mechanism is \cref{thm:irrelevance}. C4 and C1${}^{+}$ are constructed for
operational information parity, $\ladder{4}\parityeq\ladderplus$
(\cref{tab:parity-ledger}), so their information gaps vanish,
$\infogap{\ladder{4}}=\infogap{\ladderplus}=0$, and the at-optimum risks
coincide: $\bayesrisk(\ladder{4})=\bayesrisk(\ladderplus)$. At the information
limit there is simply no quantity for a format effect to be made of. What the
empirical comparison then measures is the residual of
\cref{thm:bounded-gap}: the advantage $\advantage(\ladderplus,\ladder{4})$ a
\emph{bounded} predictor can extract once the information channel is closed,
which the theorem caps at the capacity budget
$\Phi(\ladderplus,\ladder{4};\nsample,\delta)$ --- a quantity built from the
approximation gap and the induced-class Rademacher complexity, containing no
mutual information. A near-zero observed effect is the statement that this
budget was small (or unrealized) in T1's regime.

\cref{prop:regime} says when that is expected. A near-zero advantage is the
generic reading of its non-activation clause~(ii): absent a genuine capacity
separation --- a positive approximation gap $\approxerr{\ladder{4}}<
\approxerr{\ladderplus}$, or a finite-$\nsample$ estimation term comparable to
it --- the permitted window is empty and the bound itself sits near zero. A
frontier model on single-function, non-recursive Python is plausibly on the
benign side of clause~(i)'s threshold (it generalizes from in-context structure
rather than failing to), so the regime that \emph{could} have produced a large
parity-preserving effect is the one the theory marks as requiring a separation
T1 evidently did not supply. The null is thus doubly expected: forced toward
zero at the optimum by \cref{thm:irrelevance}, and not pushed away from zero by
the bounded term because T1 exhibits no capacity separation the theory can point
to.

\idealization{This subsection leans on both idealizations of
\cref{rem:slogan}. (i)~The exact zero is the \emph{Bayes-optimal} risk equality;
the frozen LLM is not the posterior $\post{\rep}$, so what is observed is the
bounded residual of \cref{thm:bounded-gap}, and the claim is only that this
residual is small, not that it is structurally zero. (ii)~Parity is operational,
so $\infogap{\cdot}=0$ holds up to the \cref{sec:parity-defect} defect rather
than exactly; the C1${}^{+}$/C4 byte asymmetry of \cref{prop:abs-strict} is the
admitted slack. The mapping of $\widehat f_{\rep}$ onto a single model's
in-context predictions is the analogy flagged in the preamble.}

\subsection{The T3 ceiling as sub-frontier saturation}
\label{subsec:anchor-t3}

T3 records a larger point estimate, $r_{\mathrm{rb}}=+0.43$, but with
$93.7\%$ of scores pinned at the maximum $1.0$ and an effective sample of
$\nsample_{\mathrm{eff}}\approx 5$ after the ceiling censors the rest. It is
tempting to read the larger $r$ as a latent structure effect that T1 lacked.
The theory says the opposite, and assigns T3 to a \emph{different} regime than
T1 --- this distinctness is the substantive content of the subsection.

T1's null lives at the information limit (\cref{thm:irrelevance}): the channel
is closed and the bounded residual is small. T3's ceiling lives instead at the
\emph{floor of the risk}, where the bounded-learner gap is inactive for a reason
unrelated to parity. When a task sits below the model's capacity frontier, the
predictor realizes (near-)Bayes behaviour in \emph{every} representation:
$\risk(\widehat f_{\rep})\to\bayesrisk(\mathrm{id})$ with
$\approxerr{\rep}\to 0$ for each C-condition simultaneously. By
\cref{cor:realizable-gap} the achievable advantage then collapses to the
estimation term alone, and with the risk already at its floor there is no
headroom in which any term --- capacity or otherwise --- can express itself: the
saturation is the \emph{realizable, infinite-resource} corner of
\cref{prop:regime}(iii), reached not by taking $\nsample\to\infty$ but by the
task being easy enough that finite resources already suffice. The bounded gap is
inactive because the frontier was never approached, which is a categorically
different statement from T1's ``inactive because the information channel is
closed.''

On top of this theoretical saturation sits a measurement one, and the two must
not be conflated. The $93.7\%$ ceiling censors the score distribution, so even
if a sub-ceiling advantage existed it could not be resolved: with
$\nsample_{\mathrm{eff}}\approx 5$ the $r_{\mathrm{rb}}=+0.43$ is an estimate
over a handful of unsaturated items and carries almost no power. The honest
reading is therefore that T3 is \emph{uninformative} about the bounded gap in
both senses at once --- the theory expects no active gap (sub-frontier
saturation) \emph{and} the instrument could not have measured one (ceiling
censoring). The larger $r$ is consistent with the theory not because the theory
predicts it but because the theory predicts the cell carries no identifying
signal in either direction; reading the point estimate as evidence of a
structure advantage would be an overclaim the theory forbids.

\idealization{The saturation argument uses \cref{cor:realizable-gap}, which
\emph{relaxes} idealization~(i) (it is a bounded-learner statement, not a Bayes
one) but assumes the base class renders $Y$ realizable in every coordinate ---
the formal content of ``below the capacity frontier.'' For a frozen LLM this
realizability is an empirical surmise about task difficulty, not a verified
property; it is the conjectural step here, and the
distinctness-from-T1 claim depends on it.}

\subsection{The T2 positive as a permitted, unidentified bounded effect}
\label{subsec:anchor-t2}

For completeness the third cell: T2 recorded the largest clean effect,
$r_{\mathrm{rb}}=+0.58$ at raw $p=0.023$, which did not survive the Holm
correction over the nine-cell confirmatory family ($m=9$; non-rejection on all
nine). The theory neither claims nor needs this as a positive. What it says is
narrow and one-directional: an effect of this magnitude is \emph{permitted} ---
it lies inside the capacity budget $\Phi$ of \cref{thm:bounded-gap} and the
two-sided envelope $|\advantage|\le\Phi$ of \cref{cor:size-gap}, so its mere
existence violates no bound --- and, were it real, it would sit in the
activation window of \cref{prop:regime}(ii), the finite-sample
capacity-separated regime the theory names as the only place a parity-preserving
effect can live. That is the whole of the retrodiction: the sign is
unconstrained by the theory, the size is bounded, and the
statistical verdict on the frozen record is non-rejection. T2 is the cell most
worth pointing a designed bounded-regime experiment at (\cref{sec:predictions}),
precisely because it is the one the theory cannot resolve from the existing
data.

\subsection{Limits of the mapping}
\label{subsec:anchor-limits}

\begin{remark}[What this section does and does not establish]
\label{rem:anchor-limits}
The three cells map cleanly onto three named regimes --- T1 to the
information-limit invariance of \cref{thm:irrelevance} (channel closed), T3 to
sub-frontier saturation via \cref{cor:realizable-gap} and
\cref{prop:regime}(iii) (no headroom, compounded by ceiling censoring), and T2
to the permitted-but-unidentified activation window of \cref{prop:regime}(ii).
The mapping's limits, stated plainly so \cref{sec:limitations} can collect them:
\begin{enumerate}
  \item[(L1)] \emph{Consistency is not confirmation.} Every claim above is of
    the form ``the record is consistent with the theory.'' A null and a ceiling
    are weak evidence; they fail to refute, which is not the same as
    corroborating. The theory earns its keep through the \emph{predictions} of
    \cref{sec:predictions}, not through this retrodiction.
  \item[(L2)] \emph{The theory does not prove any effect is zero.}
    \cref{thm:irrelevance} forces the \emph{Bayes-optimal} risks equal; it is
    silent on a bounded learner, where \cref{thm:bounded-gap} permits a nonzero
    advantage up to $\Phi$. ``No information-channel effect is expected'' is the
    exact claim; ``the effect is zero'' is not, and conflating them would be the
    central overclaim this paper takes care to prevent.
  \item[(L3)] \emph{Idealizations are not met by a frozen LLM.} The model is
    neither Bayes-optimal (idealization~(i)) nor at absolute parity
    (idealization~(ii), corrected only up to the \cref{sec:parity-defect}
    defect), and it is a single fixed predictor rather than an ERM over a class
    with a sample-size knob, so $\nsample$, $\hypclass$, and $\widehat f_{\rep}$
    enter by analogy. The retrodiction is robust to these only because it asserts
    \emph{bounds and regimes}, never point values.
  \item[(L4)] \emph{The T1/T3 distinction is partly conjectural.} That T3 sits
    below the capacity frontier (hence realizable, hence saturated) is an
    empirical surmise about task difficulty, not a theorem; if
    T3 were in fact frontier-hard, its ceiling would be measurement saturation
    without the sub-frontier reading, and only the censoring argument would
    survive.
\end{enumerate}
In one line: the theory retrodicts \emph{which regime each cell occupies} and
\emph{why no information-limit effect should appear}, and it declines ---
deliberately --- to retrodict the existence, sign, or magnitude of any
bounded-regime effect, leaving those to the falsifiable tests of
\cref{sec:predictions}.
\end{remark}

% Section 8 — Predictions and falsifiable tests (charter §4.8; issue T3-2)
\section{Predictions and falsifiable tests}
\label{sec:predictions}

% STUB (T3-2). The paper's forward-looking contribution and the antidote to
% "a math-dressed null" (charter §1 positive half; guardrail G2).
%
% Each prediction carries (charter §4.8):
%   - a named bounded-regime experiment the theory implies,
%   - a SIGN and SIZE prediction (size bounded by \cref{prop:advantage}),
%   - an explicit FALSIFIER.
%
% Hard rule (guardrail G2): these are falsifiable PREDICTIONS, never
% established empirical findings. Anything beyond the proved theorems is
% marked as conjecture (guardrail G6).

\noindent\emph{[Stub. Fills in T3-2. Predictions derive from
\cref{sec:bounded}; each is stated with sign, size bound, and falsifier.]}

% Section 9 — Related work (charter §4.9; issue T4-2, WOV-275)
\section{Related work}
\label{sec:related}

% Citation hygiene (guardrail G5, from R6, VERBATIM):
%   - Do NOT cite "SOAP" as a work.
%   - Do NOT mischaracterize INDICT (Le et al., NeurIPS 2024).
%   - Verify Clover, Plan-and-Solve, PAL before citing.
%   - Find correct references independently rather than reusing prior
%     confabulated exemplars.
% Every \cite in this section was verified against a primary source at
% insertion time (see theory.bib header). The section positions the paper; it
% does not survey. Each area is cited only where it touches a specific result,
% and no individual ingredient is claimed as this paper's novelty (guardrail
% G2); the circumscribed claim is stated in \cref{subsec:rw-novelty}.

This paper sits at the meeting point of four literatures: information-theoretic
accounts of what a representation can carry, learning-theoretic accounts of what
a learner can extract from it, the categorical machinery we borrow to state
parity exactly, and the empirical study of code representations for language
models that motivates the whole question. We situate each in turn, and in each
case the relationship is the same in form: the prior work supplies an object or a
tradeoff, and our contribution is to fix one axis (target information) so the
other (a bounded learner's value) can be isolated.

\subsection{Information-theoretic views of representation}
\label{subsec:rw-info}

The negative half of the paper (\cref{sec:parity}) is built from sufficiency and
the data-processing inequality, both standard \citep{coverthomas2006,
lehmanncasella1998}: a sufficient statistic for $Y$ retains all predictive
information, and post-processing cannot increase mutual information. The
information-bottleneck programme \citep{tishby2000informationbottleneck}, and its
variational deep-learning realisation \citep{alemi2017deepvib}, turn this calculus
into a design objective by trading compression of $X$ against retention of
information about $Y$ along a Lagrangian frontier. Our setting is the
complementary one. We do not move along that frontier; we hold the predictive
term $\MI{Y}{\rep(X)}$ \emph{fixed} across representations (information parity,
\cref{def:parity}) and ask what is left to vary. The irrelevance theorem
(\cref{thm:irrelevance}) is, in bottleneck language, the statement that two points
on the same level set of retained information are indistinguishable to the Bayes
optimum, so any residual effect cannot be an information-bottleneck effect at all.

\subsection{Learning-theoretic analyses of representation}
\label{subsec:rw-learning}

A second body of work studies, within statistical learning theory, when a
\emph{learned} representation provably helps: multitask representation learning
quantifies the sample-complexity benefit of a shared feature map
\citep{maurer2016multitask}, and transfer-learning theory shows the benefit hinges
on task diversity \citep{tripuraneni2021taskdiversity}. These results ask which
representation to \emph{learn} and prove that the right one lowers estimation
error. The positive half of this paper (\cref{sec:bounded}) asks a different
question on the same machinery: given two representations of equal information,
what advantage can a \emph{fixed}, finite-capacity learner extract from one over
the other? Our bound (\cref{thm:bounded-gap}) is in the same currency these
analyses use --- an approximation term and a Rademacher estimation term
\citep{bartlettmendelson2002, shalevshwartz2014, mohri2018} --- but it carries no
mutual-information term, which is exactly what marks the residual as a capacity
phenomenon rather than an information one.

\subsection{Applied category theory in machine learning}
\label{subsec:rw-cat}

The go/no-go section on parity (\cref{sec:category}) reads operational parity as a
retraction and tracks its round-trip defect. We use category theory only as a
language for that one structural fact, and we want the borrowing to be honest
about its weight. There is a substantial programme that does treat learning
itself categorically --- supervised learning as a functor
\citep{fong2019backpropfunctor} and categorical foundations of gradient-based
learning via lenses and reverse-derivative categories
\citep{cruttwell2022categorical}. We make no contribution to that programme. The
categorical content here is the thin observation that a section/retraction pair
exposes a defect invisible to both the mutual information and the Bayes risk
(\cref{thm:retraction}); the standard references \citep{maclane1998, riehl2016,
awodey2010} supply everything we use, and the section survives only because it
earns a theorem the core sections cannot state (charter §2).

\subsection{Code representation for language models}
\label{subsec:rw-code}

The empirical question this paper formalises is the one studied directly by its
companion preprint, whose related-work section (\texttt{preprint-v1}, §2) we
extend rather than repeat. Two lines bear on it. \emph{Training-time} integration
of code structure bakes ASTs, data flow, or tree structure into model weights ---
CodeBERT \citep{feng2020codebert}, GraphCodeBERT \citep{guo2021graphcodebert},
TreeBERT \citep{jiang2021treebert}, and the generative InCoder
\citep{fried2023incoder}. The closest controlled ancestor is
\citet{namavar2022coderepr}, who compare token, AST, and graph representations for
learning-based program repair and find the better representation is task- and
model-dependent rather than universal. We do not retrain; we vary the input
representation of a frozen model and, unlike that lineage, hold information
constant by construction, so a null is informative rather than inconclusive.

\emph{Inference-time} structure enters at the prompt boundary. Chain-of-thought
\citep{wei2022chainofthought}, plan-then-solve scaffolding
\citep{wang2023plansolve}, and program-aided prompting \citep{gao2023pal} all
improve a fixed model by changing the form, not the information content, of the
input --- and PAL in particular adds an external execution resource, which places
its gains squarely on the capacity/operational side of our account rather than the
information side. A parallel line documents that frontier models are sensitive to
semantically inert prompt-format changes \citep{sclar2024promptformatting,
salinas2024butterfly}; this is precisely the bounded-learner swap gap our positive
half predicts (\cref{sec:bounded}) and our parity-defect section bounds
(\cref{sec:parity-defect}). Finally, two neighbouring areas are sometimes
conflated with representation effects but are operational-consistency mechanisms:
contract- and verifier-backed consistency checking of generated code
\citep{sun2024clover}, and deterministic AST analysis for detecting and correcting
code hallucinations \citep{khati2026asthallucination}. We cite these for what they
are; neither bears on the information-parity claim.

\subsection{Novelty, circumscribed}
\label{subsec:rw-novelty}

Every ingredient above pre-exists: sufficiency and the DPI, the bottleneck
tradeoff, representation-learning bounds, the categorical view of learning, and
the empirical comparison of code formats. The companion preprint's contribution is
a controlled empirical \emph{combination} --- information-parity testing of a
frozen model under pre-registered correction. This paper's contribution is
narrower and theoretical: a formal account of \emph{when} code-representation
format can and cannot help a bounded language model --- never at the information
limit (\cref{thm:irrelevance}), and otherwise only within a capacity budget that
carries no information term and is non-vacuous only in a named finite-resource
regime (\cref{thm:bounded-gap,prop:regime}) --- anchored to a controlled null
rather than to a manufactured positive. We claim the localisation and its
falsifiers (\cref{sec:predictions}), not any of the components it is built from.

% Section 10 — Limitations and honesty (charter §4.10; issue T4-3)
% T4-3 is the manuscript's INTEGRITY GATE (guardrails G1/G2/G3) and must not
% be skipped (charter §6 note).
\section{Limitations and honesty}
\label{sec:limitations}

% STUB (T4-3). Assumptions and failure modes; what the theory does NOT
% establish; internal red-team against overclaim.
%
% Checklist this section must discharge:
%   - G1: the null is a prediction, not a disappointment — verify no section
%     framed it as a result the method "should have" turned positive.
%   - G2: no manufactured positive claim — verify §8 stays predictions, not
%     findings.
%   - G3: every invocation of \cref{thm:irrelevance} discloses its
%     idealization (Bayes-optimality + exact parity); collect them here.
%   - G6: every claim beyond a proved theorem is marked conjecture.
%   - Demotion record: if §3 or §6 failed go/no-go, say so plainly here.

\noindent\emph{[Stub. Fills in T4-3, the integrity gate. Red-team the whole
manuscript against overclaim before this section closes.]}

% Section 11 — Conclusion (charter §4.11; issue T4-1)
\section{Conclusion}
\label{sec:conclusion}

We set out to explain, rather than merely report, the central null of
\texttt{preprint-v1}, and to do so without converting a null into a manufactured
positive. The explanation is two-halved. The negative half is a theorem: under
information parity the data-processing inequality closes at the optimum, and the
Bayes risk is format-invariant (\cref{thm:irrelevance}), so the Lumen T1 null is
the \emph{predicted} information-limit outcome and not an anomaly to be explained
away (guardrail~G1). When parity fails, the failure is not catastrophic but
graded: the inflation of Bayes risk is the lost information exactly under log
loss and within a $\sqrt{\cdot}$ envelope otherwise (\cref{cor:ratedist-gap}).
The positive half localizes, and bounds, the only place a real effect can live.
A bounded learner is not the Bayes posterior, and the room between them is a
capacity budget with no information content (\cref{thm:bounded-gap},
\cref{cor:size-gap}); that budget is non-empty only in the finite-resource,
capacity-separated window the theory names (\cref{prop:regime}). The two go/no-go
refinements sharpen the picture without changing it: operational parity is a
retraction whose defect hides from both the mutual information and the Bayes risk
(\cref{thm:retraction}), and that same defect, once metrized, dominates the swap
gap of any fixed bounded predictor (\cref{thm:defect-risk-gap}).

Read together, the three frozen Lumen cells occupy three distinct regimes ---
T1 at the closed information channel, T3 at sub-frontier saturation, T2 in the
permitted-but-unidentified activation window (\cref{sec:anchoring}) --- and we
claim of this mapping only consistency, never confirmation
(\cref{rem:anchor-limits}). The theory earns its keep not by retrodiction but by
risk: the five predictions of \cref{sec:predictions} attach a sign-aware trend,
a prohibition, or a bound to a concrete knob of the apparatus, each with an
explicit falsifier (\cref{tab:pred-gates}). A parity-verified effect that
\emph{grows} as the model recedes from its frontier would falsify the thesis; a
clean null in the frontier regime would corroborate the negative half; a
confirmed monotone effect would scope the product to the bounded regime the
theory marks out. None of these outcomes leaves the theory ``always consistent,''
which is the point.

The honest residue is small but real, and \cref{sec:limitations} collects it.
The irrelevance theorem lives at two idealizations a frozen LLM does not meet;
the bounded-learner objects map onto a single model by analogy; the T1/T3
distinction rests on an empirical surmise about task difficulty that is
conjectural, not proved (guardrail~G6). What survives all of this is a single,
falsifiable claim: format cannot buy information that is not there, so any
genuine code-representation effect on a bounded model is a capacity phenomenon of
bounded size, appearing only near the model's frontier --- and that is a
statement the apparatus can be pointed at and made to test.


\appendix
% Appendix — Full proofs (charter §4 appendix; issue T4-4)
\section{Proofs}
\label{app:proofs}

% STUB (T4-4). Full proofs of all theorems, lemmas, and propositions, each
% with its idealization disclosure attached (guardrail G3).
%
% As theorems land in T2, their proofs migrate here and the body keeps only
% the statement. Each proof block ends with the \idealization{...} note that
% matches its theorem.
%
% Index of obligations (update as theorems are proved):
%   - \cref{thm:irrelevance} (irrelevance at optimum, T2-2)
%   - \cref{prop:advantage} (bounded-regime advantage bound, T2-4)
%   - parity-defect propositions (T2-5, if §6 passes go/no-go)
%   - retraction proposition (T2-3, if §3 passes go/no-go)

\noindent\emph{[Stub. Proofs migrate here in T4-4 as theorems are
established.]}


\bibliographystyle{plainnat}
\bibliography{theory}

\end{document}
